%%%%%%%%%%%%%%%%%%%%%%%%%%%%%%%%%%%%%%%%%%%%%%%%%%%%%%%%%%%%%%%%%%%%%%%
% !TEX TS-program = xelatex
% !TEX encoding = UTF-8
% !TEX root = aims_syllabus.tex
% For any queries, contact Papa Sega WADE (https://papasegawade.com)
%%%%%%%%%%%%%%%%%%%%%%%%%%%%%%%%%%%%%%%%%%%%%%%%%%%%%%%%%%%%%%%%%%%%%%%

\section{Course Title Placeholder}\label{course-title-placeholder}

\textbf{Institution}: AIMS Senegal \\
\textbf{Course Load}: 30 hours \\
\textbf{Target Audience}: Master II students

\begin{center}\rule{0.5\linewidth}{0.5pt}\end{center}

\subsection{1. How to use this template}\label{how-to-use}

This is a placeholder body file (\texttt{course\_body.tex}). In a real-world scenario, you do not need to write this entire layout in pure LaTeX!

You can write your syllabus easily in Markdown (like \texttt{my\_syllabus.md}) and generate this file automatically using \textbf{Pandoc}.

For example, run this in your terminal:
\begin{verbatim}
pandoc my_syllabus.md -t latex --top-level-division=section -o course_body.tex
\end{verbatim}

The \texttt{main.tex} file will then automatically include your generated content while keeping the beautiful AIMS branded cover page and headers intact!

\subsection{2. Example Table}\label{example-table}

If you have tables that look squished, you can manually adjust the \texttt{longtable} width ratios in the generated LaTeX file.

\begin{longtable}[]{@{}ll@{}}
\toprule\noalign{}
Topic & Hours \\
\midrule\noalign{}
\endhead
\bottomrule\noalign{}
\endlastfoot
Foundations & 10h \\
Advanced Topics & 20h \\
\end{longtable}
